\clearpage
\phantomsection\addcontentsline{toc}{chapter}{原课程资料与课历}
\markboth{原课程资料与课历}{}
\noindent{\color{structurecolor}\bfseries\fontsize{14.4}{20}\selectfont 原课程资料与课历}\par

\noindent\textbf{6.253: Convex Analysis and Optimization, Spring 2012.}
MIT 电气工程与计算机科学系研究生课程, 教师 Dimitri P. Bertsekas.
每周两次讲课, 每次1.5小时. 官方网页没有逐讲公历日期或上课钟点, 因而下面保留讲次顺序, 不填推测日期.
先修为一门线性代数课(最好包含抽象理论)及一门实分析课, 例如18.06、18.100.
课程强调严格证明、几何直觉及算法的适用条件. 成绩由作业25\%、期中考试25\%、学期论文50\%组成.
主教材为 Bertsekas 的 \emph{Convex Optimization Theory}, Athena Scientific, 2009.
课程大纲的阅读顺序为: 基本凸性(1.1--1.4节), 凸性与优化(第3章), 超平面与共轭(1.5--1.6节), 多面体凸性(第2章), 几何对偶(第4章), Lagrange/Fenchel/锥对偶及鞍点(5.1--5.3、5.5节), 次梯度与最优性(5.4节), 凸优化算法(第6章).
这里的教材章节号属于原教材, 与本册重构章号不同.

\begin{longtable}{@{}p{0.8cm}p{10.8cm}p{3.9cm}@{}}
\toprule 讲次 & 原课主题 & 本册对应\\\midrule\endhead
1 & 凸性在优化中的作用; 对偶; 算法与对偶 & 第\ref{chap:1}、\ref{chap:4}章\\
2 & 凸集与函数; 上图; 闭凸函数; 凸性的判别 & 第\ref{chap:1}章\\
3 & 可微凸函数; 凸包与仿射包; Carathéodory定理 & 第\ref{chap:1}章\\
4 & 相对内部与闭包; 二者的运算; 凸函数连续性与闭包 & 第\ref{chap:2}章\\
5 & 衰退锥与线性空间; 函数衰退方向; 局部/全局极小; 最优解存在 & 第\ref{chap:2}章\\
6 & 闭集交非空; 解的存在; 线性变换下闭性; 超平面 & 第\ref{chap:3}章\\
7 & 超平面分离; 非竖直超平面; 共轭函数; 共轭定理 & 第\ref{chap:3}章\\
8 & 共轭复习; 最小公共点/最大截距(MC/MC); 弱对偶 & 第\ref{chap:3}、\ref{chap:4}章\\
9 & 极小极大与零和博弈; MC/MC定理; 强对偶; 对偶解存在 & 第\ref{chap:4}章\\
10 & MC/MC定理III; 非线性Farkas引理; 线性及凸规划对偶; 最优条件 & 第\ref{chap:4}章\\
11 & 凸规划对偶反例; Fenchel对偶; 锥对偶 & 第\ref{chap:4}、\ref{chap:6}章\\
12 & 次梯度; Fenchel不等式; 灵敏度; 次微分运算; 最优性 & 第\ref{chap:5}章\\
13 & 问题结构; 锥规划 & 第\ref{chap:6}章\\
14 & 锥与半正定规划; 精确罚函数; 下降与最速下降 & 第\ref{chap:7}章\\
15 & 次梯度计算、迭代与收敛 & 第\ref{chap:7}章\\
16 & 近似次梯度; 近似方法; 切平面 & 第\ref{chap:7}、\ref{chap:8}章\\
17 & 切平面; 单纯形分解; 二者的对偶关系 & 第\ref{chap:8}章\\
18 & 广义多面体近似; 结合切平面与单纯形分解 & 第\ref{chap:8}章\\
19 & 近端极小化及其扩展 & 第\ref{chap:9}章\\
20 & 近端切平面; 束方法; 增广Lagrange; 对偶近端 & 第\ref{chap:9}章\\
21 & 广义近端; 内点与障碍法; 锥规划内点 & 第\ref{chap:9}--\ref{chap:11}章\\
22 & 大规模求和; 增量梯度/次梯度/近端; 循环与随机选择 & 第\ref{chap:10}章\\
23 & 梯度投影; 迭代复杂度; 外推; 最优一阶算法 & 第\ref{chap:11}章\\
24 & 梯度近端; 非二次近端; 熵、指数增广Lagrange与镜像下降 & 第\ref{chap:11}章\\
25 & 凸分析、对偶与算法复习 & 第\ref{chap:13}章\\\bottomrule
\end{longtable}

以上课历属于2012年讲义. \href{https://ocw.mit.edu/courses/6-253-convex-analysis-and-optimization-spring-2012/pages/assignments/}{官方作业页}明确说明五组作业来自2010年春季;
本册分别标为2010年作业、2010年期中、2012年期中. 原作业文件没有截止日期, 不将它们安插进2012年某一周.
\href{https://ocw.mit.edu/courses/6-253-convex-analysis-and-optimization-spring-2012/pages/syllabus/}{课程大纲}与
\href{https://ocw.mit.edu/courses/6-253-convex-analysis-and-optimization-spring-2012/pages/lecture-notes/}{讲次索引}是这份课历的来源.

\clearpage
\noindent\textbf{15.053: Optimization Methods in Management Science, Spring 2013.}
MIT Sloan 管理学院本科课程, 教师 James Orlin 与 Ebrahim Nasrabadi.
每周两次讲课, 每次1.5小时; 每周一次习题课, 每次1.5小时. 官方无正式先修要求.
无必购教材, 原课使用 Excel Solver. 课程按线性规划、整数规划、网络流与启发式、决策树、行为经济学组织.
本册原第\ref{chap:12}章只取已选的线性规划讲次; 增补题目按原课程库存收入, 因而还包含后续整数规划、网络与决策树题, 不将这些题伪称全部属于线性规划.

考核为书面作业10\%、电子表格作业10\%、测验2--5合计10\%、测验6--7合计10\%、两次期中各20\%、团队项目15\%、课堂参与5\%.
测验1是练习, 原权重0\%, 在能提高测验平均分时可计入; 测验2--5去掉最低分后把其余总分折算为10分.
七次测验通常20--30分钟. 期中1考线性规划, 期中2考整数规划; 没有期末考试.
六组作业独立提交, 可讨论但不可抄袭; PS1、PS2各分两组, 所以前两题合并、第3题两版均保留.
习题课自愿参加; 两次期中考前各有一次复习. 团队项目要求4--6页团队论文、1页个人学习报告及展示.
这些是2013年原课规定, 不是本册给自学者新增的评分安排.

\begin{longtable}{@{}p{0.9cm}p{8.4cm}p{6.2cm}@{}}
\toprule 课次 & 原课主题 & 原课程节点\\\midrule\endhead
L1 & 课程介绍 & \\
L2 & 线性与非线性规划建模 & \\
R1 & 习题课1 & \\
L3 & 线性规划几何与图示 & 测验1(建模); PS1截止\\
L4 & 单纯形法1 & \\
R2 & 习题课2 & \\
L5 & 单纯形法2 & 测验2(建模与LP几何); PS2截止\\
R3 & 习题课3 & \\
L6 & 灵敏度与影子价格 & \\
L7 & 双人零和/常和博弈1 & 测验3(单纯形); PS3截止\\
R4 & 习题课4 & \\
L8 & 博弈论2 & \\
L9 & 项目讨论; 松弛与人工变量 & 测验4(灵敏度); PS4截止\\
R5 & 习题课5 & \\
考试 & 期中考试1 & 官方此页没有实际试卷\\
L10 & 整数规划介绍 & 项目成员名单提前一天截止\\
L11 & 整数规划建模 & 选题及简介两天后截止\\
R6 & 习题课6 & \\
L12 & 分支定界 & \\
L13 & 割平面 & 测验5(整数建模); PS5与项目初步提案截止\\
R7 & 习题课7 & \\
L14 & 整数规划建模续 & \\
L15 & 网络1: 最短路 & PS6截止\\
R8 & 习题课8 & \\
考试 & 期中考试2 & 项目最终提案截止; 提供综合练习\\
L16 & 网络2: 网络流 & \\
L17 & 网络3: 旅行商问题 & 班级票选课堂项目展示\\
R9 & 习题课9 & \\
L18 & 决策树1 & 测验6(网络)\\
L19 & 决策树2: 信息价值 & \\
R10 & 习题课10 & \\
L20 & 行为经济学1 & 测验7(决策树); 团队论文及个人报告截止\\
L21 & 行为经济学2 & \\
L22 & 项目报告 & \\
L23 & 项目报告 & 最终反馈\\\bottomrule
\end{longtable}
课历取自\href{https://ocw.mit.edu/courses/15-053-optimization-methods-in-management-science-spring-2013/pages/syllabus/}{2013年大纲}.
L为讲课, R为习题课. 大纲给课次及相对截止节点, 没有具体上课钟点或逐讲日期.
作业PDF另记PS1截止2月12日、PS2截止2月21日、PS3截止2月28日、PS4截止3月7日、PS6截止4月11日(均2013年).
PS5标题原写“Thursday April 2th, 2013”, 日期4月2日与星期四相矛盾, 无据决定哪个正确, 故保留矛盾而不猜测改成4月4日.
PS1组别文件名与页首“Second Group”亦有不一致, 按官方资源页Group 1/2及对应内容标版本. 这些日期不反推为另一个课程的日历.
习题课1--9原由 Ebrahim Nasrabadi、Giacomo Nannicini 为2012年课程编制, 经Nannicini许可收入2013年页面, 其年份出处在题章保留.

\noindent\textbf{公开考试资料的边界.}
6.253公开两份期中卷及解答. 15.053公开期中2综合练习及解答, 没有两份实际期中试卷或七份测验试卷.
不能用复习指南反造缺失试卷. 测验2指南列九类任务: 二维可行域与冗余约束, 凸形与顶点, 指定最优点的目标构造, 松弛/剩余变量, 标准单纯形表判别, 基与基本可行解, 进入变量及最优性, 沿边移动与无界判断, 主元更新.
测验3指南列七类: 进入变量, 比值与离基/无界, 退化, 替代最优解, 唯一点/最优边/最优射线, Bland规则, Phase I.
测验6指南覆盖网络术语与Dijkstra, 最小费用流、整数性、运输与指派, 独立1与覆盖线的对偶, TSP及设施选址的构造和改进启发式.
指南3说明测验25分钟; 指南6正文的“3D-minute”是原件字符错误, 不能据此捏造一项精确时长.
